\documentclass[10pt,a4paper]{amsart}
\usepackage[margin=19mm]{geometry}
\usepackage{amsmath,amssymb,mathtools}
\usepackage[scr=rsfs,cal=euler]{mathalfa}
\usepackage{xfrac}
\usepackage[unicode,hidelinks,bookmarks=false]{hyperref}
\newcommand{\ie}{i.e.\ }
\newcommand{\cf}{cf.\ }
\newcommand{\resp}{respectively\ }
\newcommand{\Subsection}{\subsection}
\providecommand{\nobreakdash}{\nobreak\hbox{-}}
\setlength{\emergencystretch}{2em}
\multlinegap0pt
\usepackage{bm}
\usepackage{bbm}
\usepackage{dsfont}
\usepackage{xspace}
\usepackage{cancel}
\usepackage{mathtools}
\usepackage{amsthm}
\makeatletter
\@ifundefined{theorem}{\newtheorem{theorem}{Theorem}}{}
\@ifundefined{lemma}{\newtheorem{lemma}[theorem]{Lemma}}{}
\makeatother
\def\F{\mathbb{F}}
\title[Three-dimensional Hadamard matrices of Paley type]{Three-dimensional Hadamard matrices of Paley type}
\author{Vedran Kr\v{c}adinac, Mario Osvin Pav\v{c}evi\'c, Kristijan Tabak}
\date{Source: arXiv:2305.12415v2 (CC BY 4.0)}
\begin{document}
\maketitle

\noindent\emph{Source and licence.} Three-dimensional Hadamard matrices of Paley type. Version arXiv:2305.12415v2 (CC BY 4.0), distributed under the Creative Commons Attribution 4.0 International licence, as recorded in the arXiv licence metadata for this exact version and shown on the abstract page https://arxiv.org/abs/2305.12415v2. Full rights notice: licenses/arxiv-2305.12415-CC-BY-4.0.txt.\par\smallskip

\noindent\textit{Editorial scope.} Counted unit: the Theorem whose source label is \texttt{maintm} in the article (source0.tex lines 213--217, source label \texttt{maintm}), delivered together with the article's own complete original proof of it (source0.tex lines 219--243, one \texttt{proof} environment of 25 lines). No mathematics is written, repaired or re-derived: statement and proof are the article's own text line for line. Required context retained uncounted: paley (lines 144--151); maindef (lines 169--180); lm1 (lines 184--187); lm2 (lines 193--197). Every definition, display and label of those lines is retained unchanged. Omitted from this package and not redistributed: 1--143, 152--168, 181--183, 188--192, 198--212, 218--218, 244--306. Nothing inside a retained excerpt is omitted or altered: every retained body is delivered exactly as the article's lines are written, so no omission receipt of any kind accompanies this package. Citations: the retained excerpts carry no citation command at all, so no bibliography entry of the article is transcribed and no reference is redistributed. Reference adaptation: none was needed and none was made; every reference inside the delivered unit resolves inside it, so no unresolved reference or citation is produced. Printed slips kept literal: no printed word, symbol, hypothesis or number of the counted statement or of its proof was altered. Layout and class: the article's own class is \texttt{amsart} with packages including amsmath, amssymb, amsfonts, amsthm, color, hyperref, graphicx, amsaddr; the wrapper is the lane's pinned \texttt{amsart} editorial wrapper at 10pt on A4, which restates the theorem-like environments the retained text needs and adds the article's own macro definitions that the retained text needs (\texttt{\textbackslash F}), with extra wrapper packages (bm, bbm, dsfont, xspace, cancel, mathtools). The delivered numbering is the unit's own. Assets: no retained span contains a figure environment, a graphics-inclusion command, a PDF-page inclusion, a table or a TikZ picture; no image file is copied into this package, the package declares no asset dependency and no image file is redistributed with it. Licence: version arXiv:2305.12415v2 of this article is distributed under the Creative Commons Attribution 4.0 International licence, as recorded in the arXiv licence metadata for this exact version and shown on the abstract page https://arxiv.org/abs/2305.12415v2. A full rights notice is delivered at licenses/arxiv-2305.12415-CC-BY-4.0.txt. Characters: every retained excerpt is pure ASCII, so the delivered engine, XeLaTeX with its default Latin Modern text font, reports no missing character.
\begin{equation}\label{paley}
h(x,y)=\left\{\begin{array}{ll}
-1, & \mbox{if } x=y=\infty,\\
1, & \mbox{if } x=y\neq \infty \mbox{ or } x=\infty\neq y \\
 & \mbox{ or } y=\infty\neq x,\\
\chi(y-x), & \mbox{otherwise.}
\end{array}\right.
\end{equation}

\begin{equation}\label{maindef}
H(x,y,z)=\left\{\begin{array}{ll}
-1, & \mbox{if } x=y=z,\\
1, & \mbox{if } x=y\neq z\\
 & \mbox{ or } x=z\neq y\\
 & \mbox{ or } y=z\neq x,\\
\chi(z-y), & \mbox{if } x=\infty,\\
\chi(x-z), & \mbox{if } y=\infty,\\
\chi(y-x), & \mbox{if } z=\infty,\\
\chi((x-y)(y-z)(z-x)), & \mbox{otherwise.}
\end{array}\right.
\end{equation}

\begin{lemma}\label{lm1}
The matrix~\eqref{maindef} is invariant under cyclic shifts of the coordinates,
i.e.\ $H(x,y,z)=H(y,z,x)=H(z,x,y)$.
\end{lemma}

\begin{lemma}\label{lm2}
The matrix~\eqref{maindef} is invariant under linear fractional transformations
of the coordinates with determinant~$1$, i.e.\ under the action of $PSL(2,q)$ on
the projective line.
\end{lemma}

\begin{theorem}\label{maintm}
For every odd prime power~$q$, equation~\eqref{maindef} defines a three-dimensional Hadamard
matrix of order~$q+1$. If $q\equiv 3 \pmod{4}$, the matrix is proper with all $2$-dimensional
layers equivalent to the Paley type I matrix~\eqref{paley}.
\end{theorem}

\begin{proof}
Because of Lemma~\ref{lm1} we may fix the last coordinate. We claim that
for all distinct $a,b\in PG(1,q)$,
$$\sum_{x,y\in PG(1,q)} H(x,y,a) H(x,y,b)=0.$$
Now, because $PSL(2,q)$ acts $2$-transitively on $PG(1,q)$ and Lemma~\ref{lm2},
we may take $a=0$ and $b=\infty$ without loss of generality. We divide the sum
into two parts:
$$\sum_{x} H(x,x,0)H(x,x,\infty) + \sum_{x\neq y} H(x,y,0)H(x,y,\infty).$$
The first part is readily seen to be $q-3$. For the second part we distinguish whether
$0$, $\infty$ are among $x$, $y$ or not. If $x=0$, $y=\infty$ or $x=\infty$, $y=0$
the summand is $1$, so up to now the total is $q-1$. If $x\in \F_q^*$, $y=\infty$, the
sum is $$\sum_{x\in \F_q^*} H(x,\infty,0)H(x,\infty,\infty)=\sum_{x\in \F_q^*} \chi(x)=0$$
because there are equally many squares and non-squares in $\F_q^*$. Similarly we see that
the cases $x\in \F_q^*$, $y=0$; $x=\infty$, $y \in \F_q^*$; and $x=0$, $y \in \F_q^*$ sum
up to $0$. The final part of the sum is over $x,y\in \F_q^*$, $x\neq y$:
\begin{multline*}
\sum H(x,y,0)H(x,y,\infty)=\sum \chi((x-y)(y-0)(0-x))\chi(y-x)= \\[2mm]
= \sum \chi(x)\chi(y) =\sum_x \chi(x)\sum_{y\neq x}\chi(y) = \sum_x \chi(x)(-\chi(x))=1-q.
\end{multline*}
The grand total is $q-1+1-q=0$ and the first part of the theorem is proved. For the second
part notice that thanks to Lemmas~\ref{lm1} and~\ref{lm2} we may, without loss of generality,
fix the third coordinate to $z=\infty$ and look at the $2$-dimensional layer obtained by
varying $x$ and $y$. In this case equation~\eqref{maindef} reduces to~\eqref{paley}, and
this is a Paley type~I Hadamard matrix if $q\equiv 3 \pmod{4}$.
\end{proof}

\end{document}
